% ============================================================================
\clearpage
\section*{v4.9.0-r70 Source-Only: The 4.10 Consolidation in the Monolith, a Correction to r69, and Two Tests of a Common Dynamics}
\addcontentsline{toc}{section}{v4.9.0-r70 Source-Only: The 4.10 Consolidation in the Monolith, a Correction to r69, and Two Tests of a Common Dynamics}
\label{sec:v490-r70-consolidation}
% ============================================================================

\paragraph{Purpose.}
On 29 and 30 October 2026 the author continued OP5 with another language model.  Besides new working notes in \texttt{analysis/}, that work
changed the research candidate on GitHub directly (commits \texttt{f98b29b} to \texttt{cf48496}): it corrected the three main-body errors found in
v4.9.0-r69, added further status notes and a corrected Hopf annotation, and added the section ``States, Mediators and Clock Comparisons''
(\S\ref{sec:v410-state-mediator}) with a figure.  This layer carries those changes into the monolith, records them, corrects one number of
v4.9.0-r69, and records two checks of the common-dynamics sketch of 28 October: the pre-registered test H48 and an exploration of the idea that
\(\Lambda\) is a mixture over sector crossings.  Earlier layers are not rewritten.

\noindent\textbf{Source-only status.}
This v4.9.0-r70 layer is source-only.  It does not build a PDF, does not assign a new DOI,
and does not authorize a public release.  The v4.9.0 Version DOI
\href{https://doi.org/10.5281/zenodo.23534954}{10.5281/zenodo.23534954} and the Concept DOI
\href{https://doi.org/10.5281/zenodo.19210728}{10.5281/zenodo.19210728} are unchanged by it.

\subsection*{v4.9.0-r70 results}
\label{subsec:v490-r70-results}

\paragraph{Record v4.9.0-r70.1 (the 4.10 consolidation).}
The following changes were made on GitHub without a layer and are applied here verbatim (patch specification
\texttt{llm/revision\_spec\_v410\_konsolidierung.py}, eight patches; local insertion \texttt{llm/einbau\_v410\_abschnitte\_lokal.py}).
\begin{enumerate}[leftmargin=2em]
\item[(a)] Remark~\ref{rem:regularization-B}: series \(R(\tau)=-1/\tau+1-\tfrac{11}{24}\tau+O(\tau^2)\), counterterm \(-\log(\varepsilon/\mu)\), closed
finite part, value \(1.686953464\) at \(T=0.393\), \(\mu=1\); the old target of Proposition~\ref{prop:theta-required-B} is withdrawn and
Remark~\ref{rem:op5-structural-summary} is rewritten.  This writes candidate (a) of v4.9.0-r69.3.
\item[(b)] Status notes at Theorem~\ref{thm:dirac-probe} and Proposition~\ref{prop:three-compatible-readings} (archived GUE clock scale) and at the
corollary ``Cosmological constant revisited'' of \S\ref{sec:phase-constants}.  This writes candidates (b) and (c) of v4.9.0-r69.3.
\item[(c)] Further status notes at Theorem~\ref{thm:entanglement-shared-time} (a shared clock does not by itself prove entanglement) and
Corollary~\ref{cor:falsifiable} (historical prediction register), and a dependency note after the Theta-sum conjecture.
\item[(d)] Theorem~\ref{thm:hopf}: the fibre \(S^7\) was annotated as a group containing \(G_2\supset SU(3)\).  \(S^7\) is not that group; \(G_2\) is the
automorphism group of the octonions and \(SU(3)\) the stabiliser of an imaginary unit.  The annotation is removed and the reason given in a note.
\item[(e)] Section~\ref{sec:v410-state-mediator} and Subsection~\ref{subsec:op5-multimode-slices} with Figure~\ref{fig:op5-multimode-slices}.  On GitHub they
are separate parts of the candidate; in the monolith they stand before the archive appendix.  The figure file \texttt{figures/op5-slices-9x9.pdf} is the
first external graphic of the monolith.
\end{enumerate}
The section records conditional results: typed readouts and the obligations of a mediator, an orientation-complete local reconstruction of the
real doublet, calibrated clock comparisons, and finite field-compatibility tests up to fourth jet order at three points.  None is a regional
solution, a native common action or an OP5 closure; the section says so itself.

\paragraph{Correction v4.9.0-r70.2.}
The value \(1.689\) in v4.9.0-r69.3(a) is wrong; the finite part is \(1.686953464\).  The quadrature in \texttt{llm/k106\_nachrechnung\_op5.py}
integrated \(1+R+1/\tau\) near \(\tau=0\), where two large terms cancel, and lost accuracy.  The direct limit and the closed antiderivative agree
to twelve digits.  The finding of r69.3(a), that the printed expansion, counterterm and value \(0.940\) were wrong, stands.

\paragraph{Theoremlet v4.9.0-r70.3 (H48: \(\bsym\) as a frozen length).}
Pre-registered in \texttt{llm/Vorab\_Notiz\_H48\_Einfrieren\_beta.md} before computing.  Model: the M/T doublet of a homogeneous spinor condensate
is moved by the duality rotation, the compensator measurement and exchange with the tachyon reference state, with rates proportional to the density;
flat unimodular FLRW with the Einstein--Cartan contact term.
\begin{enumerate}[leftmargin=2em]
\item[(a)] Refuted.  Over 288 runs the frozen length \(|b_\infty|\) ranges from \(10^{-12}\) to \(0.87\) and depends on the initial expansion rate and density,
not only on the rate ratio \(r\).  The dense-phase value \(r/(1+r)\) is reached only if freezing happens while the expansion is still faster than the
rotation (21 runs).
\item[(b)] The Einstein--Cartan contact term changes \(|b_\infty|\) by at most \(10^{-4}\): in the dense phase the state lies on the R/T axis, and the
axial density along the compensator (chirality) axis vanishes.
\item[(c)] \(\Lambda\) changes \(|b_\infty|\) by up to 18\,\% when freezing falls into the \(\Lambda\) phase; it does not decouple from the frozen state.
\end{enumerate}

\paragraph{Remark v4.9.0-r70.4 (\(\Lambda\) over sector crossings; exploration, not a test).}
The author asked whether each crossing of sectors carries its own \(\Lambda\), so that ours is a mixture.  A linear dependence \(\mathrm{Tr}(\rho\hat\Lambda)\)
averages out under the rotation; a bilinear one depends on \(|b|\) only and is constant once \(|b|\) is frozen, so the epoch question of v4.9.0-r67 stays.
Of 26 readings with structural weights (\texttt{llm/k108\_kreuzungen.py}) one lands within \(2\sigma\) of DESI DR2 BAO
(\(\Omega_M=1/\gL^2=0.288\), \(3.9\sigma\) from Planck 2018), about as many as chance gives (\(1.3\)).  No prediction is made.  An identity remains:
with \(|b|=1/\gL\) the eigenvalues of the M/T state satisfy \(\lambda_-/\lambda_+=(\gL-1)/(\gL+1)=\bsym^2\), so \(\bsym^2\) is the weight ratio of the
mixture.  This is a rewording, not a new result.

\paragraph{Proof sketch.}
r70.1 by the patch specifications (reproducing the GitHub fragments exactly); r70.2 by direct evaluation of the finite part; r70.3 by
\texttt{llm/k107\_test\_h48.py} (output \texttt{llm/k107\_out.txt}, result \texttt{llm/Ergebnis\_H48\_Einfrieren\_beta.md}); r70.4 by
\texttt{llm/k108\_kreuzungen.py}.  \(\square\)

\subsection*{v4.9.0-r70 status in the three frames}
\label{subsec:v490-r70-frames}

\begin{itemize}[leftmargin=2em]
\item r70.1 and r70.2 concern R and the main body; the old ``photon frame'' of the corrected passages is R (v4.9.0-r24).
\item r70.3: the doublet carries M (state \(b\)) and T (\(b\mapsto-b\)); the frozen length \(|b|\) is the same in M and T, so (a) is frame-free.  The
photon P is the attractor \(b=0\).
\item r70.4: \(\lambda_\pm\) are the weights of M and T; the identity is symmetric in M and T.
\end{itemize}
No statement of this layer concerns colour or charge.

\subsection*{v4.9.0-r70 register entries}
\label{subsec:v490-r70-register}

\begin{center}
\small
\begin{tabular}{@{}>{\raggedright\arraybackslash}p{0.34\textwidth}>{\raggedright\arraybackslash}p{0.58\textwidth}@{}}
\toprule
statement & status after v4.9.0-r70\\
\midrule
status-note candidates of v4.9.0-r69.3 & written (4.10 consolidation)\\
value in v4.9.0-r69.3(a) & corrected to \(1.686953464\)\\
H48 (\(\bsym\) as frozen length) & refuted; depends on expansion history\\
\(\Lambda\) over sector crossings & no reading better than \(\Omega_M=\bsym^2\); epoch question open\\
monolith and GitHub candidate & main-body text identical again\\
\bottomrule
\end{tabular}
\end{center}

\subsection*{v4.9.0-r70 ledger and guard update}
\label{subsec:v490-r70-ledger-guard}

\[
R_{\rm chk}^{v490\text{-}r70}=(1,1,0)
\]
for (consolidation carried over and r69 corrected; H48 refuted and the crossing reading without a hit beyond chance; OP5 not closed).

The layer spends no guard:
\[
R_{\rm glob}^{(v490\text{-}r70,+)}=(0,1,1,1,1),
\qquad
\Gamma_{\rm guard}^{v490\text{-}r70}=(1,1,1,1,1),
\qquad
\beta_4=1/256\ \text{unspent}.
\]
It does not imply exact finite \(C5\) zero closure, exact all-readout terminality,
\(C6/T2\), a global photon-clock/RH theorem, or global all-readout closure.
